\chapter{Measure Spaces}

This chapter lays the measure-theoretic foundation required to build the Itô integral. 
While standard texts like Le Gall focus heavily on pure mathematical rigor, our goal here is dual : 
to preserve that formal architecture while bridging the gap toward quantitative intuition. 
Before we can price a derivative or simulate a stochastic process in Python, we must first rigorously define the space of possible events and how we measure them.

\begin{quote}
\itshape
``The starting idea of measure theory is to assign a non-negative real number (the measure of this subset) to each subset of a given set, in a way that satisfies certain natural additivity properties (the measure of a disjoint union must be the sum of the measures). 

For deep reasons, it is generally not possible to define the measure of any arbitrary subset, and one must restrict to a certain class of subsets, called \textbf{measurable sets}; a set equipped with a $\sigma$-algebra is called a \textbf{measurable space}.''
\hfill --- J.-F.~Le Gall~\cite{legall2006}
\end{quote}

% ========================================================================================================================================================================================================================================================================


% ========================================================================================================================================================================================================================================================================

\section{Measurable Sets and $\sigma$-Algebras}

\begin{definition}{$\sigma$-Algebra}{sigma_algebra}
Let $E$ be a set. 
A collection $\mathcal{A}$ of subsets of $E$ is called a \textbf{$\sigma$-algebra} on $E$ if:
\begin{enumerate}[label=(\textbf{\roman*})]
    \item \textbf{Certain Event :} $E \in \mathcal{A}$ ;
    \item \textbf{Closed under Complementation :} $A \in \mathcal{A} \implies A^c \in \mathcal{A}$ ;
    \item \textbf{Closed under Countable Unions :} If $(A_n)_{n \in \mathbb{N}} \subseteq \mathcal{A}$, then $\bigcup_{n \in \mathbb{N}} A_n \in \mathcal{A}$.
\end{enumerate}
\end{definition}

Integration relies on the passage to the limit.  
We need to sum an infinite number of small pieces to perform limit calculations.
The elements of $\mathcal{A}$ are called measurable sets, or $\mathcal{A}$-measurable sets if there is any ambiguity.
We say that $(E, \mathcal{A})$ is a \textbf{measurable space}.

\begin{insight}[$\sigma$-algebra definition]
\end{insight}

\begin{insight}[Intuition: $\sigma$-Algebras as Information Resolution]
Think of a $\sigma$-algebra $\mathcal{F}$ as the \textbf{resolution of an image} representing the market:
\begin{itemize}
    \item \textbf{Coarse resolution (Little Information):} If $\mathcal{F} = \{\emptyset, \Omega\}$, the image is completely pixelated into a single block. You only know that "something will happen" ($\Omega$), but you cannot distinguish any specific outcome.
    \item \textbf{Finer resolution (More Information):} As price observations arrive, the image sharpens. The set $\Omega$ is partitioned into smaller subsets. An event $A \in \mathcal{F}$ means your current resolution is sharp enough to tell whether the true state of the world lies inside $A$ or not.
\end{itemize}
In financial modeling, a \textbf{filtration} $(\mathcal{F}_t)_{t \ge 0}$ is simply a movie where the picture becomes clearer as time $t$ progresses.
\end{insight}

\begin{insight}[Intuition: $\sigma$-Algebras as Decidable Payoffs]
Mathematically, an event $A \subseteq \Omega$ belongs to $\mathcal{F}$ if and only if, given our current information, we can answer the question with certainty: \emph{"Did event $A$ occur?"}

In quantitative finance, if a derivative contract pays $\$1$ when $A$ occurs and $\$0$ otherwise, the contract can only be settled at time $t$ if $A \in \mathcal{F}_t$. If $A \notin \mathcal{F}_t$, the market does not yet possess enough information to determine the payoff.
\end{insight}

\section{Integration \& Limit Theorems}

\begin{theorem}{Lebesgue's Dominated Convergence Theorem}{dct}
Let $(X_n)_{n \ge 0}$ be a sequence of random variables on $(\Omega, \mathscr{F}, \mathbb{P})$. Assume that:
\begin{itemize}
    \item $X_n \to X$ almost surely as $n \to \infty$.
    \item There exists a non-negative random variable $Y \in L^1(\mathbb{P})$ such that $|X_n| \le Y$ a.s. for all $n \ge 0$.
\end{itemize}
Then $X \in L^1(\mathbb{P})$ and:
\[
\lim_{n \to \infty} \mathbb{E}[X_n] = \mathbb{E}[X]
\]
\end{theorem}

\begin{proof}
The proof relies on applying Fatou's Lemma to the non-negative sequence $Y + X_n \ge 0$ and $Y - X_n \ge 0$.
\end{proof}


\begin{example}{Borel $\sigma$-Algebra}{borel_algebra}
Let $\Omega = \mathbb{R}$ with the usual Euclidean topology. The Borel $\sigma$-algebra $\mathscr{B}(\mathbb{R})$ is the $\sigma$-algebra generated by all open sets of $\mathbb{R}$.
\end{example}

\begin{insight}[Why Borel sets?]
We want to measure probabilities of intervals (like a stock price falling between 100 and 110). The Borel $\sigma$-algebra is simply the mathematical safety net that ensures we don't accidentally create unmeasurable "monster" sets when taking unions and intersections of these intervals.
\end{insight}